\IEEEPARstart{N}{ews} analytics pipelines assign articles to standard topics early, so that the more expensive stages only process the topics that matter. The taxonomy most used by news agencies is Media Topic of the International Press Telecommunications Council (IPTC). Its top level has 17 topics, such as \emph{politics}, \emph{economy, business and finance}, and \emph{sport}~\cite{iptc2024mediatopics,kuzman2025llm}.

Kuzman and Ljube\v{s}i\'c~\cite{kuzman2025llm} released the first open multilingual IPTC classifier. They labeled articles in four languages with GPT-4o and used these labels to fine-tune XLM-RoBERTa-large~\cite{conneau2020xlmr}. The resulting model matches GPT-4o on a human-annotated test set. A follow-up study by the same group concluded that fine-tuned encoders are more practical than large language models (LLMs) for large-scale annotation~\cite{kuzman2025southslavic}. An encoder with \val{teacher.params}~million parameters is still expensive without a graphics processing unit (GPU). On one central processing unit (CPU) core the model needs \val{eff.teacher.latency_ms_p50_1c}~ms per article, and its fastest CPU configuration, an Open Neural Network Exchange (ONNX) export quantized to 8-bit integers (int8), still needs \val{eff.teacher_onnx_int8.latency_ms_p50_1c}~ms (Section~\ref{sec:res-eff}). Neither study reports this inference cost, although inference accounts for a large share of the computing demand and energy use of machine learning in production, for example 60\% of the machine learning energy use reported by Google~\cite{luccioni2024power}.

The target is a topic classifier that covers many languages, runs in milliseconds per article on an ordinary CPU, and has a small footprint and few dependencies. Existing options fall short of this. Small distilled transformers~\cite{sanh2019distilbert,wang2021minilmv2} gain about one order of magnitude. In the CPU benchmark of this study, MiniLM-L6 is \val{eff.minilm_kd.speedup_onnx} times faster than the ONNX int8 teacher. Int8 quantization alone gives a factor of \val{eff.teacher_onnx_int8.speedup_pt} over PyTorch fp32. fastText~\cite{joulin2017bag} is much faster, but it learns its embeddings from scratch and therefore does not transfer across languages. This matters here because labeled IPTC data exists for only four languages.

This paper proposes StaticKD, which distills~\cite{hinton2015distilling} the XLM-R-large IPTC classifier into a student that, after training, is a single table of 17 logits per token. A teacher that works across languages can label unlabeled news in any language at no annotation cost. Multilingual static embeddings distilled from sentence encoders~\cite{tulkens2024model2vec,wada2025static,chen2024bgem3} place translation equivalents close to each other, and looking them up costs almost nothing. News topics are also strongly lexical, which keeps bag-of-words representations competitive for topic classification~\cite{galke2022bow}. Because the classification head of StaticKD is linear, the averaged embedding and the head can be folded into a token table. Inference then reduces to tokenization and averaging table rows, which NumPy handles.

A multilingual claim requires a multilingual evaluation with labels that do not come from the teacher. The human test set of Kuzman and Ljube\v{s}i\'c is not public and covers four languages. SIB-200~\cite{adelani2024sib200} covers 205 languages, but its items are short sentences with a non-IPTC taxonomy. To fill this gap, a unified test set of \val{ts.docs} articles in \val{ts.langs} languages is built. It combines CC-NEWS articles from 2024 to 2026 that carry their newsroom section in schema.org metadata~\cite{guha2016schema} with public datasets labeled by humans or by publisher categories~\cite{adelani2023masakhanews,petukhova2023mnds,mirashi2023indicnews,hossain2025banglanews}. The labels are editorial decisions of publishers or annotators, translated onto the official IPTC tree. No model that reads the article text is involved.

This study addresses four research questions (RQs).
\begin{itemize}
  \item \textbf{RQ1 (cost):} What is the CPU inference cost of the IPTC teacher, and how much does StaticKD reduce it?
  \item \textbf{RQ2 (accuracy):} How much of the teacher accuracy does StaticKD retain on labels that do not come from the teacher, compared with baselines of similar cost, namely fastText and a support vector machine (SVM) on term frequency--inverse document frequency (TF-IDF) features, and with baselines of higher cost, namely MiniLM-L6 and DistilmBERT?
  \item \textbf{RQ3 (components):} Which distillation components determine the accuracy of this very-low-capacity student: soft labels, multilingual data, embedding initialization, embedding fine-tuning, token dropout, temperature, or ensembling?
  \item \textbf{RQ4 (generalization):} On which languages, labels, and label sources is StaticKD weak, and which factors explain it?
\end{itemize}

Distilling into a student of such low capacity also shows where its accuracy comes from. The knowledge that matters comes from the coverage of the transfer set and from the aligned initialization, not from the shape of the soft labels. Hard labels match soft labels, and a high temperature hurts. Transfer to languages that never appear in the distillation data is carried by embedding rows that training never updates, which works because they share a space with the trained rows. Fine-tuning the embeddings trades this transfer for in-distribution fit, a representation drift that a development set in the training languages cannot detect.

The contributions of this paper are as follows.
\begin{enumerate}
  \item A multilingual IPTC test set of \val{ts.docs} articles in \val{ts.langs} languages whose editorial labels do not come from the teacher. It comes with a section-to-IPTC mapping table that is validated automatically and reviewed by the first author, cross-source deduplication, a contamination check against all training data, and site-cluster bootstrap confidence intervals, and it can be reused to evaluate other IPTC classifiers.
  \item The findings above, supported by ablations with five seeds per variant and by analyses by language, label, label origin, and source, including languages that are absent from the distillation data.
  \item StaticKD, a recipe that distills a multilingual transformer teacher into a token-logit table initialized from cross-lingual static embeddings. The final model takes \val{final.size}~MB, needs only \texttt{numpy} and \texttt{tokenizers}, and processes \val{eff.statickd.docs_per_sec_1c} articles per second per core. An ensemble of seeds folds into the same table at no extra inference cost. Its accuracy--cost trade-off is compared with nine configurations from six model families, measured on CPU only with core pinning for a hybrid CPU.
  \item A release of the code, a notebook that reproduces every number in this paper, the models, and the test-set manifest.
\end{enumerate}

Section~\ref{sec:related} reviews related work, Section~\ref{sec:method} describes the data, test set, method, and evaluation protocol, Section~\ref{sec:results} presents the results, Section~\ref{sec:discussion} discusses them, Section~\ref{sec:limitations} lists the limitations, and Section~\ref{sec:conclusion} concludes.
